\section{Introduction}
\label{sec:introduction}

Personalized pricing is a common decision problem in settings where a seller offers configurable products to buyers with heterogeneous preferences. A buyer may value the same bundle differently depending on budget, use case, feature priorities, and willingness to negotiate. This makes pricing more than a one-shot demand-estimation problem: the seller must decide how aggressively to price while observing only partial evidence about the buyer's latent preferences.

This report studies personalized pricing through a controlled benchmark prototype. The concrete substrate is customized vehicle-feature pricing, where a seller negotiates over a fixed feature bundle. The domain is useful because it combines feature-level value, implementation cost, buyer heterogeneity, and multi-round negotiation. However, the intended benchmark question is broader than vehicle customization: how should pricing agents be evaluated when buyer preferences are hidden and can only be inferred through interaction?

PrefBench is positioned between dynamic pricing, negotiation environments, and agent benchmark design. It provides a simulator-based benchmark for seller-side pricing decisions in a structured negotiation setting, with fixed assets, fixed persona splits, and shared evaluation metrics.

The contribution of this report is intentionally scoped. We release and document a working benchmark prototype rather than claiming a complete benchmark suite. The prototype includes a product catalog, persona generation pipeline, frozen persona bank, multi-round negotiation environment, and baseline evaluation protocol. We also outline how the benchmark can be extended to LLM agents, belief probing, robustness evaluation, and human reference studies.

The main contributions are:
\begin{itemize}
  \item We define a simulator-based benchmark prototype for personalized pricing and negotiation under hidden buyer preferences.
  \item We implement a reproducible environment with a structured catalog, semi-synthetic buyer personas, fixed data splits, and shared episode-level metrics.
  \item We report checkpoint-free heuristic baseline results under a common evaluation protocol and define the next LLM-facing evaluation path.
  \item We document the limitations of the current prototype and identify concrete future extensions for LLM-facing evaluation, belief probing, robustness testing, and human reference comparison.
\end{itemize}
